% Chapter 2

\chapter{Introducción específica}
\label{Chapter2}

En el presente capítulo se describen los recursos de software, los conjuntos de datos y la infraestructura de cómputo de terceros que sustentan el sistema documentado en esta memoria. Se justifica la elección de cada componente y se caracteriza el origen de la telemetría utilizada para el entrenamiento.

\section{Frameworks y bibliotecas de IA/ML}
\label{sec:frameworks-bibliotecas}

La implementación del autoencoder LSTM se apoyó en \texttt{TensorFlow} y su interfaz \texttt{Keras} \citep{abadi2016tensorflow, chollet2015keras}, que ofrecen una capa de alto nivel para definir, entrenar y serializar redes neuronales profundas. Entre las capas utilizadas figuran \texttt{LSTM}, \texttt{RepeatVector} y \texttt{TimeDistributed}, componentes nativos que permiten expresar la arquitectura codificador--decodificador con pocas líneas. La elección del \textit{framework} obedeció a la integración estable con el ecosistema Python usado para procesamiento de datos, a la portabilidad del modelo entrenado a contenedores en CPU y al formato de serialización \texttt{.weights.h5}, que se acompaña de un archivo JSON con la definición estructural del grafo para reconstruir el modelo en inferencia sin acoplarse a una versión particular del entorno de entrenamiento.

Para el preprocesamiento se utilizó \texttt{scikit-learn} \citep{pedregosa2011scikit}. La normalización se realizó con \texttt{MinMaxScaler}, configurado con cotas fijas para desacoplar la transformación de la distribución particular de los datos de entrenamiento. La persistencia del preprocesador, necesaria para garantizar la misma transformación entre entrenamiento e inferencia, se gestionó con \texttt{joblib} \citep{joblib_site}, biblioteca de serialización binaria habitual en ese ecosistema.

La manipulación de las series temporales se apoyó en \texttt{NumPy} \citep{harris2020numpy} y \texttt{pandas} \citep{mckinney2010pandas}, estándares en Python para arreglos n-dimensionales y tablas indexadas por tiempo. La obtención de métricas desde Prometheus se realizó vía HTTP con la biblioteca \texttt{requests} \citep{requests_site}, mientras que la instrumentación del servicio simulado utilizó el cliente oficial \texttt{prometheus\_client} \citep{prometheus_client_site}, que expone contadores e histogramas en el formato esperado por el recolector. La carga de archivos de configuración en formato YAML se delegó en \texttt{PyYAML} \citep{pyyaml_site}, y la generación de gráficos de evaluación \textit{offline} se realizó con \texttt{matplotlib} y \texttt{seaborn} \citep{matplotlib_site, seaborn_site}. 

En entornos de observabilidad, la elección del stack analítico no es independiente del stack de despliegue: bibliotecas con dependencias pesadas o con requisitos de GPU dificultan la ejecución periódica en contenedores pequeños. Se priorizaron componentes maduros en CPU, serialización explícita de preprocesadores y modelos, y clientes HTTP livianos para consultar Prometheus sin introducir capas intermedias.La tabla~\ref{tab:bibliotecas} resume las bibliotecas principales y el rol que cumplen dentro del sistema.

\begin{table}[h]
\centering
\small
\caption{Bibliotecas de terceros utilizadas en el sistema y su rol.}
\label{tab:bibliotecas}
\begin{tabular}{@{}p{0.26\textwidth}cp{0.46\textwidth}@{}}
\toprule
\textbf{Biblioteca} & \textbf{Versión} & \textbf{Rol en el sistema} \\
\midrule
\texttt{TensorFlow} / \texttt{Keras} & 2.8 & Definición y entrenamiento del autoencoder LSTM. \\
\addlinespace
\texttt{scikit-learn} & 1.0 & Normalización (\texttt{MinMaxScaler}) con cotas fijas. \\
\texttt{joblib} & 1.1 & Persistencia del preprocesador entrenado. \\
\addlinespace
\texttt{NumPy} & 1.19 & Operaciones con arreglos n-dimensionales y cálculo del error de reconstrucción. \\
\texttt{pandas} & 1.3 & Manipulación de series temporales y construcción de ventanas. \\
\addlinespace
\texttt{requests} & 2.26 & Consultas HTTP a la API de Prometheus. \\
\texttt{prometheus\_\allowbreak client} & 0.17 & Exposición de métricas en el servicio simulado. \\
\addlinespace
\texttt{PyYAML} & 5.4 & Carga de archivos de configuración. \\
\addlinespace
\texttt{matplotlib} / \texttt{seaborn} & 3.4 / 0.11 & Gráficos de evaluación offline del modelo. \\
\bottomrule
\end{tabular}
\end{table}

En la práctica, las dependencias se agruparon en tres bloques funcionales y modulares: cómputo numérico y aprendizaje profundo (\texttt{TensorFlow}, \texttt{NumPy}, \texttt{pandas}, \texttt{scikit-learn}), observabilidad (\texttt{requests}, \texttt{prometheus\_client}, \texttt{PyYAML}) y evaluación visual (\texttt{matplotlib}, \texttt{seaborn}). Esta separación facilitó auditar qué componentes intervienen en el entrenamiento, en la inferencia y en los ensayos documentados en el capítulo de resultados.

\section{Datasets y datos de entrenamiento}
\label{sec:datasets}

El conjunto de datos utilizado para entrenar y evaluar el sistema corresponde a telemetría operativa del servicio de TV-over-IP de la empresa cliente, recolectada con un intervalo de muestreo de 30~segundos sobre una ventana histórica de 90~días. La fuente primaria es la instancia interna de Prometheus de la organización \citep{brazil2018prometheus, prometheus_site}. Las series se obtuvieron mediante la API HTTP de rango \texttt{api/v1/query\_range}, con las consultas PromQL definidas en el archivo \texttt{config/data.yaml} del repositorio.

\subsection{Confidencialidad y entorno de demostración}

El entrenamiento sobre telemetría real se realizó dentro de la infraestructura autorizada de la empresa. No obstante, las métricas de desempeño operativo (carga, latencia, errores y consumo de recursos) constituyen información confidencial: su publicación en un repositorio académico o en capturas de la memoria permitiría inferir aspectos del comportamiento productivo del servicio.

Por ese motivo, el material reproducible del trabajo, incluido el repositorio entregable y las figuras ilustrativas de este capítulo, emplea un entorno demostrativo con un servicio simulado y un generador sintético que replica la misma familia de métricas y la semántica de agregación PromQL del entorno productivo.

Esta separación no implica que el modelo se haya entrenado únicamente con datos sintéticos: el régimen nominal utilizado para el aprendizaje provino de Prometheus interno. El simulador permitió desarrollar y validar el \textit{pipeline} de forma reproducible, ejecutar demostraciones sin acceso a la red de producción y documentar el sistema sin exponer series identificables para los fines académicos de este trabajo.

\subsection{Datasets de entrenamiento utilizados}

Al tratarse de telemetría propia de un servicio en producción, no se utilizaron \textit{datasets} públicos como fuente principal. El dominio (calidad de servicio en \textit{streaming} sobre infraestructura específica) no encuentra equivalencia directa en repositorios abiertos como Numenta NAB \citep{numenta_nab} o Yahoo Webscope \citep{yahoo_webscope}, que agrupan series heterogéneas con distinto muestreo y sin el stack Prometheus del cliente.

\section{Modelos preentrenados y transferencia de aprendizaje}
\label{sec:transfer-learning}

El sistema no incorporó modelos preentrenados ni esquemas de \textit{transfer learning}. La razón es triple. En primer lugar, no existen modelos públicamente disponibles entrenados sobre telemetría operativa de servicios de TV-over-IP con la misma combinación de métricas y régimen de muestreo de 30~segundos. La transferencia desde dominios distintos (texto, visión o series financieras) no resultaba directamente aplicable porque las entradas, la escala temporal y la semántica de las variables difieren sustancialmente.

En segundo lugar, \textit{benchmarks} genéricos como Numenta NAB entrenan o evalúan sobre series univariadas o multivariadas heterogéneas que no coinciden con el vector de once variables construido en este trabajo (cinco métricas de servicio y seis atributos temporales). Un modelo preentrenado en esos conjuntos no habría capturado las correlaciones entre latencia, tasa de errores y patrones circadianos del servicio concreto.

En tercer lugar, el volumen disponible de datos del régimen nominal resultó adecuado para un entrenamiento desde cero en tiempos razonables sobre hardware modesto, sin necesidad de inicializaciones de gran capacidad. Ante la ausencia de un punto de partida compatible y la suficiencia de datos propios, se entrenó el autoencoder LSTM exclusivamente sobre telemetría del entorno autorizado.

\section{Infraestructura de cómputo}
\label{sec:infraestructura}

El sistema fue concebido para ejecutarse en infraestructura propia, sin envío de telemetría a servicios externos. Tanto el entrenamiento como la inferencia se realizaron en CPU. El tamaño reducido del modelo y de la ventana de inferencia hicieron innecesario el uso de aceleradores gráficos.

El entorno de desarrollo se orquestó con Docker y Docker Compose \citep{merkel2014docker}, lo que permitió desplegar de forma reproducible los servicios involucrados desde un único archivo de declaración. Se distinguieron dos perfiles de uso: el perfil \texttt{dev}, con el stack completo de demostración (servicio simulado, Prometheus, Grafana y detector), y el despliegue productivo sobre la infraestructura AWS de la empresa cliente, donde únicamente se agregó el contenedor de detección a la red interna existente.

El detector se distribuyó como una imagen basada en \path{python:3.10-slim}, con un usuario sin privilegios y montajes de los directorios de modelos, configuración y registros para preservar artefactos entre reinicios. La asignación de recursos definida en \path{docker-compose.yml} contempló un límite de una CPU virtual y 2~GB de memoria RAM para el contenedor de detección, lo que resultó suficiente para una inferencia cada 30~segundos sobre ventanas de 10~minutos. Los servicios auxiliares (servicio simulado, Prometheus y Grafana) se aislaron en una red puente dedicada y se expusieron solo los puertos necesarios para inspección manual durante el desarrollo.

El entrenamiento se ejecutó en un equipo de desarrollo con procesador de uso general y sin GPU dedicada. El ciclo completo se completó en un tiempo del orden de minutos a decenas de minutos, según la cantidad de ventanas resultantes del período de 90~días. Esta capacidad de reentrenar en un horizonte corto fue un criterio explícito de diseño, dado que el régimen de operación del servicio puede evolucionar con cambios de versión del producto o con variaciones estacionales del tráfico.

La integración con Opsgenie para la emisión de alertas y la generación de enlaces contextuales a Grafana se realizó mediante sus APIs públicas a través de HTTPS, sin agentes adicionales en el contenedor. De este modo, la infraestructura agregada por el sistema se redujo a un único servicio adicional dentro de la red interna, lo que respetó los requerimientos de seguridad y privacidad establecidos para el trabajo.

La tabla~\ref{tab:infraestructura} resume el entorno técnico bajo el que se entrenó y desplegó el sistema.

\begin{table}[h]
\centering
\small
\caption{Entorno de cómputo utilizado para entrenamiento y despliegue.}
\label{tab:infraestructura}
\begin{tabular}{@{}p{0.28\textwidth}p{0.64\textwidth}@{}}
\toprule
\textbf{Componente} & \textbf{Descripción} \\
\midrule
Lenguaje y entorno base & Python 3.10 sobre imagen \path{python:3.10-slim}. \\
Orquestación de servicios & Docker Compose con perfil \texttt{dev} para entorno completo de pruebas. \\
Recursos del contenedor de detección & 1.0 CPU virtual y 2~GB de memoria RAM (límite en \path{docker-compose.yml}). \\
Aceleración por GPU & No utilizada. Todo el cómputo se realizó en CPU. \\
Servicios auxiliares & Prometheus, Grafana y servicio simulado para pruebas. \\
Red & Red puente \texttt{monitoring} en Docker, sin exposición pública. \\
Destino de despliegue productivo & Infraestructura propia de la empresa cliente sobre AWS, sin dependencia de servicios analíticos de terceros. \\
\bottomrule
\end{tabular}
\end{table}

\section{Métricas de evaluación}
\label{sec:metricas-evaluacion}

La evaluación del sistema combinó dos planos complementarios. En el primero se valoró la calidad de reconstrucción del autoencoder, base del aprendizaje no supervisado, mediante el error cuadrático medio (MSE) como función de pérdida y señal escalar de anomalía, y el error absoluto medio (MAE) como métrica auxiliar de monitoreo. En el segundo se valoró la capacidad de detección sobre ventanas con anomalías sintéticas inyectadas, para lo cual se emplearon métricas de clasificación binaria habituales: precisión, \textit{recall}, \textit{F1-score} y tasa de falsos positivos. La distinción entre ambos planos respondió al carácter no supervisado del entrenamiento: las métricas de clasificación solo pudieron computarse en escenarios controlados con anomalías conocidas, mientras que el error de reconstrucción se midió de forma continua sobre tramos nominales.

La tabla~\ref{tab:metricas-evaluacion} resume las métricas utilizadas, su definición y el propósito que cumplieron en la validación del sistema.

\begin{table}[h]
\centering
\small
\caption{Métricas de evaluación utilizadas para la validación del sistema.}
\label{tab:metricas-evaluacion}
\begin{tabular}{@{}p{0.24\textwidth}p{0.30\textwidth}p{0.36\textwidth}@{}}
\toprule
\textbf{Métrica} & \textbf{Definición} & \textbf{Propósito en la evaluación} \\
\midrule
MSE & Error cuadrático medio entre la ventana de entrada y su reconstrucción. & Función de pérdida del entrenamiento y señal escalar de anomalía en inferencia. \\
\addlinespace
MAE & Error absoluto medio entre entrada y reconstrucción. & Métrica auxiliar de monitoreo durante el entrenamiento, menos sensible a valores atípicos. \\
\addlinespace
Precisión & Fracción de ventanas marcadas como anómalas que efectivamente lo eran. & Controlar la proporción de falsas alertas elevadas al servicio de guardia. \\
\addlinespace
\textit{recall} & Fracción de anomalías reales correctamente detectadas. & Medir la cobertura de incidentes ante anomalías sintéticas inyectadas. \\
\addlinespace
\textit{F1-score} & Media armónica de precisión y exhaustividad. & Resumir en un único valor el compromiso entre falsas alertas y cobertura. \\
\addlinespace
Tasa de falsos positivos & Fracción de ventanas nominales marcadas erróneamente como anómalas. & Verificar la estabilidad del umbral sobre tramos de régimen nominal. \\
\bottomrule
\end{tabular}
\end{table}

El procedimiento concreto de cálculo de estas métricas, la calibración del umbral sobre la distribución de errores de validación y los resultados cuantitativos obtenidos se detallan en los capítulos de diseño y de resultados.


